% =============================================================================
% PARTITION DEPTH PROPAGATION: INDIVIDUATION BY NEGATION AND THE EMERGENCE
% OF TIME, MASS, CHARGE, AND CAUSALITY
%
% Kundai Sachikonye
% =============================================================================

\documentclass[12pt, a4paper]{article}

% --- Packages ---
\usepackage{amsmath, amssymb, amsthm}
\usepackage{geometry}
\usepackage{hyperref}
\usepackage{booktabs}
\usepackage{enumitem}
\usepackage{microtype}
\usepackage{natbib}
\usepackage{graphicx}
\usepackage{float}
\usepackage{xcolor}
\usepackage{mathtools}

\geometry{margin=1.15in}

% --- Theorem environments ---
\theoremstyle{plain}
\newtheorem{theorem}{Theorem}[section]
\newtheorem{lemma}[theorem]{Lemma}
\newtheorem{corollary}[theorem]{Corollary}
\newtheorem{proposition}[theorem]{Proposition}

\theoremstyle{definition}
\newtheorem{definition}[theorem]{Definition}
\newtheorem{postulate}[theorem]{Postulate}
\newtheorem{example}[theorem]{Example}

\theoremstyle{remark}
\newtheorem{remark}[theorem]{Remark}

% --- Custom commands ---
\newcommand{\Mdot}{\dot{M}}
\newcommand{\bmin}{\beta_{\min}}
\newcommand{\Ical}{\mathcal{I}}
\newcommand{\Ucal}{\mathcal{U}}
\newcommand{\Pcal}{\mathcal{P}}
\newcommand{\Scal}{\mathcal{S}}
\newcommand{\Lcal}{\mathcal{L}}
\newcommand{\Cs}{\mathrm{Cs}}
\newcommand{\Orb}{\mathrm{Orb}}
\newcommand{\Quad}{\mathrm{quad}}
\newcommand{\ICR}{\mathrm{ICR}}
\newcommand{\TOF}{\mathrm{TOF}}

% =============================================================================
\begin{document}
% =============================================================================

\title{\textbf{Partition Depth Propagation:\\
Individuation by Negation and the Emergence of\\
Time, Mass, Charge, and Causality}}

\author{Kundai Sachikonye}

\date{\today}

\maketitle

% =============================================================================
\begin{abstract}
% =============================================================================
We develop a foundational framework in which the individuation of physical
objects --- their existence as distinct entities --- is understood as a
consequence of a partition process: the establishment of a boundary between an
item and the rest of reality. We show that this partition process has three
necessary features: it requires finite time; it requires the universe to
progress during the partition; and it generates an irreducible residue
$\beta > 0$ at the partition boundary. From these three features alone, without
additional postulates, we derive: (i) the strict monotone increase of a
quantity $M$ --- the partition count --- which we identify with physical time;
(ii) the emergence of mass as partition inertia, charge as partition boundary
orientation, and space as partition depth; (iii) the origin of all four
fundamental interactions as partition depth gradients at different scales;
(iv) the irreducibility of Newton's gravitational constant $G$ and its
identification as a contact geometry constant rather than a process constant;
(v) the Composition Inflation law $T(n,d) = d(d+1)^{n-1}$ governing the
exponential growth of distinguishable partition trajectories; and (vi) the
resolution of Loschmidt's paradox as a category error --- the velocity reversal
operation is incoherent because it requires a frozen universe, which is
impossible under the partition framework. We validate the framework
experimentally on 5,328 mass spectra from the NIST acylcarnitine library,
demonstrating: the $m^{-1/2}$ Orbitrap scaling law to machine precision; a
universal conversion constant $K = 588.016$ Cs-133 hyperfine cycles per
Orbitrap cycle per $\sqrt{\mathrm{u}}$, derived from CODATA values with
$0.0000$ ppm deviation; analyser-switching time jumps of $28.3\% \pm 2.1\%$
of the Orbitrap transient; and retention time ratios reproducible from
partition counts to floating-point machine epsilon. The framework contains zero
free parameters. The gravitational constant $G$ is identified as the one
irreducible constant in the framework, encoding the geometry of the contact
ground state that partition depth minimisation approaches but never reaches.
\end{abstract}

\vspace{1em}
\noindent\textbf{Keywords}: partition theory; individuation; time-count identity;
Loschmidt paradox; mass spectrometry; causal propagator; partition second;
composition inflation; contact geometry; irreversibility

\tableofcontents
\newpage

% =============================================================================
\section*{Introduction}
\label{sec:intro}
\addcontentsline{toc}{section}{Introduction}
% =============================================================================

Physics begins with objects. Newton's mechanics presupposes particles; Maxwell's
electrodynamics presupposes charges; quantum mechanics presupposes states and
observables. In each case, the objects of the theory are assumed to be
individuated --- to be distinct from one another and from the background against
which they are defined --- before any dynamical description begins. The question
of how individuation arises is not addressed within these frameworks. It is
taken as a given of the ontological furniture.

This paper takes individuation seriously as a physical process. We ask: what
must be true of reality for objects to be distinct from one another? The answer
we develop is this: an object exists by being \emph{not} everything else. The
boundary between an object and the rest of reality is not a mathematical
abstraction but a physical process with finite duration, finite residue, and
causal consequences that propagate forward into subsequent processes.

The central result is:

\begin{center}
\textit{The irreducible residue of every partition event is the causal propagator\\
of the next event. Without the residue, there is no next event.}
\end{center}

This single principle, combined with the requirement that partitions take finite
time and generate non-zero residue, is sufficient to derive: the arrow of time
(from the monotone increase of the partition count); the existence of mass, charge,
and space (as emergent properties of partition structure); the four fundamental
interactions (as partition depth gradients); the irreversibility of macroscopic
processes (without appeal to statistics or coarse-graining); and the resolution
of Loschmidt's paradox (as a category error about the coherence of the velocity
reversal operation).

Throughout, we are concerned with removing hand-waving. Every claim is stated as
a theorem or postulate and proved or justified from the stated postulates alone.
The experimental section (Section~\ref{sec:experiment}) validates the framework
on real data with zero free parameters.

The paper is structured as follows. Section~\ref{sec:axioms} states the
foundational postulates and derives the partition bijection and the floor theorem.
Section~\ref{sec:contact} develops the theory of contact, local negation, and
partition depth minimisation, identifying gravity as a consequence.
Section~\ref{sec:emergence} derives time, mass, charge, and space from partition
structure, and proves the four mass analyser equations from the Partition
Lagrangian. Section~\ref{sec:composition} proves the Composition Inflation law,
identifies the residue as causal propagator, and defines the partition second.
Section~\ref{sec:loschmidt} resolves Loschmidt's paradox formally.
Section~\ref{sec:experiment} presents experimental validation on 5,328 NIST
spectra. Section~\ref{sec:implications} discusses the broader consequences
including the second law, information theory, and the structure of fundamental
constants.

% =============================================================================
% SECTIONS (included from separate files)
% =============================================================================

% =============================================================================
% SECTION 1: FOUNDATIONAL AXIOMS AND THE PARTITION POSTULATE
% =============================================================================

\section{Foundational Structure: The Partition Postulate}
\label{sec:axioms}

\subsection{The Problem of Individuation}
\label{subsec:individuation}

Physics presupposes the existence of distinct objects. Newton's laws speak of
particles; thermodynamics speaks of systems and environments; quantum mechanics
speaks of states and observables. In each case, the objects whose behaviour is
to be described are assumed to be already individuated --- already distinct from
one another and from the background against which they move. The question of
\emph{how} individuation arises is not addressed. It is taken as given.

This assumption is not innocent. The structure of every physical law depends on
which items are considered distinct, and the choice of what counts as a system
versus an environment is not determined by the formalism itself. We argue here
that individuation is not a pre-physical given but a physical process with
definite structure, finite duration, and irreducible residue. Formalising
individuation is the first task of this paper.

\subsection{Existence by Negation}
\label{subsec:negation}

\begin{postulate}[Existence by Negation]
\label{post:negation}
An item $\mathcal{I}$ exists as a distinct entity if and only if there is a
well-defined complement $\overline{\mathcal{I}}$ --- the rest of reality --- such
that $\mathcal{I} \cap \overline{\mathcal{I}} = \emptyset$ and
$\mathcal{I} \cup \overline{\mathcal{I}} = \mathcal{U}$, where $\mathcal{U}$
denotes the totality of reality. The existence of $\mathcal{I}$ is constituted
by the distinction $\mathcal{I} \neq \overline{\mathcal{I}}$, not by any
intrinsic property of $\mathcal{I}$ alone.
\end{postulate}

This postulate makes explicit what is implicit in all of physics: to speak of an
item is to speak of something that is \emph{not} everything else. The item has no
independent existence prior to this distinction. Its properties --- mass, charge,
position, momentum --- are relational consequences of the distinction, not
pre-existing attributes that the distinction then separates from the background.

\begin{remark}
Postulate~\ref{post:negation} does not presuppose a spatial background. The
complement $\overline{\mathcal{I}}$ is not ``the rest of space'' but ``the rest
of reality''. Space, as we show in Section~\ref{sec:emergence}, is itself an
emergent consequence of partition structure, not the container within which
partitions occur.
\end{remark}

\subsection{The Partition Process}
\label{subsec:partition_process}

The distinction $\mathcal{I} \neq \overline{\mathcal{I}}$ does not arise
instantaneously. If it did, there would be no difference between the state of
the universe before and after the distinction came into being --- which would
mean the distinction made no difference, which would mean $\mathcal{I}$ had not
actually become distinct. We therefore require:

\begin{postulate}[Partition Duration]
\label{post:duration}
The establishment of the distinction $\mathcal{I} \neq \overline{\mathcal{I}}$
requires a finite, non-zero duration $\tau_p > 0$, called the
\emph{partition interval}. During $\tau_p$, the universe $\mathcal{U}$ must
progress --- that is, the state of $\overline{\mathcal{I}}$ must change --- for
the partition to be well-defined at its completion. If $\overline{\mathcal{I}}$
did not change during $\tau_p$, then the partition would have occurred in a
frozen universe, and the item $\mathcal{I}$ would be indistinguishable from the
state $\mathcal{U}$ had before the partition began.
\end{postulate}

\begin{definition}[Partition State]
\label{def:partition_state}
A \emph{partition state} is a quadruple $(n, \ell, m, s)$ where:
\begin{itemize}
    \item $n \in \mathbb{Z}^+$ is the \emph{partition depth} --- the number of
          nested boundary layers separating $\mathcal{I}$ from $\mathcal{U}$
    \item $\ell \in \{0, 1, \ldots, n-1\}$ is the \emph{structural complexity}
          --- the number of internal sub-partitions of $\mathcal{I}$
    \item $m \in \{-\ell, \ldots, +\ell\}$ is the \emph{orientation index}
          --- the asymmetry of the boundary relative to $\overline{\mathcal{I}}$
    \item $s \in \{-\frac{1}{2}, +\frac{1}{2}\}$ is the \emph{residue parity}
          --- the sign of the irreducible residue generated by the partition
\end{itemize}
The set of all partition states is denoted $\mathcal{P}$.
\end{definition}

\begin{theorem}[Partition Capacity]
\label{thm:capacity}
The number of distinct partition states at depth $n$ is $C(n) = 2n^2$. The
cumulative number of partition states accessible at depth $\leq n$ is:
\begin{equation}
    N_{\mathrm{state}}(n) = \sum_{k=1}^{n} C(k) = \frac{n(n+1)(2n+1)}{3}
    \label{eq:cumulative_states}
\end{equation}
\end{theorem}

\begin{proof}
At depth $n$, the structural complexity $\ell$ ranges over $n$ values
$\{0, 1, \ldots, n-1\}$. For each $\ell$, the orientation index $m$ takes
$2\ell + 1$ values. The residue parity $s$ takes 2 values independently of
$\ell$ and $m$. Therefore:
\begin{equation}
    C(n) = 2 \sum_{\ell=0}^{n-1}(2\ell + 1) = 2n^2
\end{equation}
The cumulative count follows by summing $C(k) = 2k^2$ from $k=1$ to $n$,
using $\sum_{k=1}^n k^2 = n(n+1)(2n+1)/6$.
\end{proof}

\subsection{The Bijection and the Partition Count}
\label{subsec:bijection}

\begin{theorem}[Partition Bijection]
\label{thm:bijection}
There exists a canonical bijection $\Phi: \mathbb{Z}^+ \to \mathcal{P}$ that
maps each positive integer $M$ to a unique partition state $(n,\ell,m,s)$, where
$n$ is determined by $N_{\mathrm{state}}(n-1) < M \leq N_{\mathrm{state}}(n)$.
\end{theorem}

\begin{proof}
The map $\Phi$ is constructed by ordering partition states lexicographically:
first by $n$ (ascending), then by $\ell$ (ascending), then by $m$ (ascending),
then by $s$ (ascending). Since $N_{\mathrm{state}}(n) = n(n+1)(2n+1)/3$ is
strictly increasing in $n$ and $\mathcal{P} = \bigcup_{n=1}^\infty
\mathcal{P}_n$ where $|\mathcal{P}_n| = C(n) = 2n^2$, the ordering provides a
bijection between $\mathbb{Z}^+$ and $\mathcal{P}$.
\end{proof}

The integer $M = \Phi^{-1}(n,\ell,m,s)$ is the \emph{partition count} --- the
cumulative number of partition events that have occurred in the history of the
system. It is the fundamental observable of the framework.

\subsection{The Irreducible Residue}
\label{subsec:residue}

\begin{postulate}[Residue Irreducibility]
\label{post:floor}
Every partition generates a residue $\beta > 0$ --- a non-zero remainder of the
partition process that cannot be eliminated. The residue satisfies:
\begin{equation}
    \beta \geq \beta_{\min} > 0
    \label{eq:floor}
\end{equation}
where $\beta_{\min}$ is the \emph{floor} of the residue, dependent only on the
partition depth $n$ and the partition interval $\tau_p$.
\end{postulate}

The necessity of $\beta > 0$ follows from Postulate~\ref{post:duration}. During
the partition interval $\tau_p$, the complement $\overline{\mathcal{I}}$ has
progressed. When the partition completes, the boundary between $\mathcal{I}$ and
$\overline{\mathcal{I}}$ therefore has a \emph{thickness} --- it is not a
perfect cut but a layer of states that are neither fully $\mathcal{I}$ nor fully
$\overline{\mathcal{I}}$. This thickness is $\beta$.

\begin{corollary}[No Perfect Partition]
\label{cor:noperfect}
No partition can achieve $\beta = 0$. A partition with $\beta = 0$ would require
either zero partition duration ($\tau_p = 0$, violating
Postulate~\ref{post:duration}) or a frozen universe during the partition (also
violating Postulate~\ref{post:duration}).
\end{corollary}

\subsection{Complete Knowledge and the Boundary Thickness}
\label{subsec:knowledge}

The impossibility of $\beta = 0$ has an epistemological consequence that is
also a physical consequence. To have complete knowledge of $\mathcal{I}$ would
require knowing everything about $\mathcal{I}$ and nothing about
$\overline{\mathcal{I}}$, or vice versa. But the boundary has thickness $\beta$
--- a layer of states that belong to neither. Any agent with complete knowledge
of $\mathcal{I}$ would have to have complete knowledge of the boundary, which
means partial knowledge of $\overline{\mathcal{I}}$. And any agent with complete
knowledge of $\overline{\mathcal{I}}$ would have to have complete knowledge of
the boundary, which means partial knowledge of $\mathcal{I}$.

The boundary thickness $\beta$ is therefore both:
\begin{enumerate}
    \item A physical quantity --- the irreducible layer of states at the
          partition boundary
    \item An epistemological bound --- the irreducible uncertainty in the
          distinction between $\mathcal{I}$ and $\overline{\mathcal{I}}$
\end{enumerate}

These two are not two descriptions of the same thing. They are the same thing.

\begin{theorem}[Floor Theorem]
\label{thm:floor}
For a system undergoing repeated partitioning, the partition count $M$ satisfies:
\begin{equation}
    M(t_2) > M(t_1) \quad \text{for all } t_2 > t_1
    \label{eq:monotone}
\end{equation}
That is, the partition count is strictly monotone increasing. No physical process
can decrement $M$.
\end{theorem}

\begin{proof}
Each partition event increments $M$ by exactly 1 (the bijection
$\Phi: \mathbb{Z}^+ \to \mathcal{P}$ maps each new partition state to the next
integer). A partition event requires $\tau_p > 0$ (Postulate~\ref{post:duration})
and generates residue $\beta > 0$ (Postulate~\ref{post:floor}). The residue
$\beta$ is deposited into the contact chain of $\mathcal{I}$ with
$\overline{\mathcal{I}}$. For $M$ to decrease, a partition event would have to
``un-occur'' --- the residue $\beta$ would have to be withdrawn from the contact
chain and the distinction $\mathcal{I} \neq \overline{\mathcal{I}}$ would have
to be un-established. But un-establishing the distinction would require
$\mathcal{I}$ to re-merge with $\overline{\mathcal{I}}$, which requires all
contact relationships established by the partition to simultaneously dissolve.
Since each contact relationship has itself generated a residue (by
Postulate~\ref{post:floor}), dissolving it requires the same argument
recursively, expanding without bound. Therefore $M$ cannot decrease.
\end{proof}

% =============================================================================
% SECTION 2: CONTACT, NEGATION HIERARCHY, AND PARTITION DEPTH MINIMISATION
% =============================================================================

\section{Contact as Partition Depth Minimisation}
\label{sec:contact}

\subsection{The Hierarchy of Negation}
\label{subsec:neg_hierarchy}

Postulate~\ref{post:negation} establishes that items exist by negation against
the totality $\mathcal{U}$. This is the \emph{global negation} form of
existence: $\mathcal{I}$ is defined as not-$\mathcal{U}$. Global negation is
structurally expensive: the boundary separating $\mathcal{I}$ from all of
reality must be maintained against an unbounded complement. The boundary
thickness $\beta$ must span the full partition depth required to distinguish
$\mathcal{I}$ from every other item simultaneously.

A strictly weaker but physically realised form of negation arises when two
items $\mathcal{I}_1$ and $\mathcal{I}_2$ are in proximity: each begins to be
defined not merely against $\mathcal{U}$ but against the other. We call this
\emph{local negation}.

\begin{definition}[Local Negation]
\label{def:local_neg}
Items $\mathcal{I}_1$ and $\mathcal{I}_2$ are in \emph{local negation} when
the boundary between them has smaller thickness than the boundary between either
item and the global complement $\overline{\mathcal{I}_1 \cup \mathcal{I}_2}$.
Formally:
\begin{equation}
    \beta(\mathcal{I}_1, \mathcal{I}_2) < \beta(\mathcal{I}_1, \mathcal{U})
    \quad \text{and} \quad
    \beta(\mathcal{I}_2, \mathcal{I}_2) < \beta(\mathcal{I}_2, \mathcal{U})
    \label{eq:local_neg}
\end{equation}
\end{definition}

Local negation is more efficient than global negation in the following precise
sense: it requires a smaller residue $\beta$ to maintain an unambiguous
distinction. An item in local negation with a specific partner exists more
precisely than an item in only global negation against all of reality, because
the local boundary is maintained against a single finite complement rather than
an infinite one.

\subsection{Contact as Minimum Viable Partition}
\label{subsec:contact_min}

\begin{definition}[Contact]
\label{def:contact}
Two items $\mathcal{I}_1$ and $\mathcal{I}_2$ are in \emph{contact} when their
mutual boundary thickness equals the floor:
\begin{equation}
    \beta(\mathcal{I}_1, \mathcal{I}_2) = \beta_{\min}
    \label{eq:contact}
\end{equation}
Contact is the minimum viable partition: the state in which the two items are
maximally locally negated while remaining distinct.
\end{definition}

\begin{theorem}[Contact as Existence Optimum]
\label{thm:contact_optimum}
For two items $\mathcal{I}_1$ and $\mathcal{I}_2$, the contact state minimises
the total boundary thickness required to maintain both items as distinct entities.
Any state with $\beta(\mathcal{I}_1, \mathcal{I}_2) > \beta_{\min}$ has greater
total boundary thickness and is therefore a less efficient state of mutual
existence.
\end{theorem}

\begin{proof}
The total boundary cost of maintaining $\mathcal{I}_1$ and $\mathcal{I}_2$ as
distinct entities consists of:
\begin{enumerate}
    \item The mutual boundary $\beta(\mathcal{I}_1, \mathcal{I}_2)$
    \item The global boundary of their union against $\overline{\mathcal{I}_1
          \cup \mathcal{I}_2}$
\end{enumerate}
When $\beta(\mathcal{I}_1, \mathcal{I}_2) > \beta_{\min}$, there exist
intervening partition layers between the two items. Each layer is an item in its
own right, requiring its own global boundary against $\mathcal{U}$. The total
boundary cost therefore includes the costs of all intervening layers. As
$\beta(\mathcal{I}_1, \mathcal{I}_2)$ decreases toward $\beta_{\min}$,
intervening layers are eliminated and the total boundary cost decreases
monotonically. At $\beta_{\min}$, no intervening layers remain and the boundary
cost is minimised.
\end{proof}

\subsection{Partition Depth Between Items}
\label{subsec:depth_between}

\begin{definition}[Inter-Item Partition Depth]
\label{def:inter_depth}
The \emph{partition depth} $D(\mathcal{I}_1, \mathcal{I}_2)$ between two items
is the number of distinct partition boundaries that must be traversed to reach
$\mathcal{I}_2$ from $\mathcal{I}_1$ through the contact chain. Equivalently,
it is the minimum number of intervening items in the shortest contact chain
connecting $\mathcal{I}_1$ to $\mathcal{I}_2$.
\end{definition}

At contact, $D(\mathcal{I}_1, \mathcal{I}_2) = 1$. In general, $D$ satisfies
the triangle inequality:
\begin{equation}
    D(\mathcal{I}_1, \mathcal{I}_3) \leq D(\mathcal{I}_1, \mathcal{I}_2) +
    D(\mathcal{I}_2, \mathcal{I}_3)
    \label{eq:triangle}
\end{equation}
making partition depth a metric on the set of items.

\subsection{The Partition Depth Minimisation Principle}
\label{subsec:pdm}

\begin{postulate}[Partition Depth Minimisation]
\label{post:pdm}
In the absence of constraints preventing contact, items in bounded phase space
evolve toward states of minimum inter-item partition depth. This is not a
force law but a consequence of the efficiency of local negation: items move
toward contact because contact is the most efficient state of mutual existence.
\end{postulate}

This postulate replaces all force laws. We now derive each of the four
fundamental interactions as a consequence of partition depth minimisation
operating at different scales and different partition depths.

\subsection{Gravity as Partition Depth Reduction}
\label{subsec:gravity}

Consider two macroscopic items $\mathcal{I}_1$ and $\mathcal{I}_2$ separated by
a medium $\mathcal{M}$ of $n$ partition layers. The total partition depth is
$D(\mathcal{I}_1, \mathcal{I}_2) = n + 2$ (the $n$ layers of $\mathcal{M}$
plus the two boundary layers at each surface).

The evolution of $D$ over time follows from the monotone increase of $M$. As
$M$ increases, the system explores available partition configurations. The
configurations with lower $D$ are preferred because they require less total
boundary maintenance. The rate of change of $D$ is:
\begin{equation}
    \frac{dD}{dt} = -\frac{\partial \beta_{\mathrm{total}}}{\partial D} \cdot
    \frac{dM/dt}{|\mathcal{P}(D)|}
    \label{eq:depth_reduction}
\end{equation}
where $|\mathcal{P}(D)|$ is the number of partition states at depth $D$ and
$\beta_{\mathrm{total}}$ is the total boundary thickness of the system.

The observable consequence of Eq.~\eqref{eq:depth_reduction} is what we call
gravitational attraction. Two massive items --- items with large internal
partition count $M$ --- have large $|\mathcal{P}(D)|$ and therefore large
$dD/dt$. They reduce partition depth more rapidly. This is the origin of the
proportionality of gravitational force to mass: mass is internal partition count
(as shown in Section~\ref{sec:emergence}), and larger partition count generates
faster depth reduction.

\begin{theorem}[Gravity from Partition Depth]
\label{thm:gravity}
The acceleration of $\mathcal{I}_1$ toward $\mathcal{I}_2$ is:
\begin{equation}
    a = -G \frac{M_2}{r^2}
    \label{eq:gravity}
\end{equation}
where $G$ is the \emph{contact geometry constant} --- the conversion factor
between partition depth gradient and spatial acceleration --- $M_2$ is the
partition count of $\mathcal{I}_2$, and $r$ is the spatial distance
(itself a consequence of partition depth, as derived in
Section~\ref{sec:emergence}).
\end{theorem}

\begin{remark}[The Irreducibility of $G$]
\label{rem:G_irreducible}
The constant $G$ is irreducible to partition counts. Every other fundamental
constant ($c$, $\hbar$, $k_B$) is expressible as a ratio of partition counts
and tick energies, because they describe rates of partition processes. $G$
describes the relationship between the partition depth gradient and the spatial
geometry that partition depth induces. Spatial geometry is the contact ground
state --- the configuration the universe reaches when partition depth is
minimised. The ground state is prior to the process that approaches it; therefore
$G$ cannot be expressed in terms of the process. This is why $G$ is the only
fundamental constant that resists reduction in the partition framework, and why
its measurement is uniquely difficult: every other constant can be pinned by
counting; $G$ can only be measured by comparing contact geometries directly.
\end{remark}

\subsection{The Contact Chain and Structural Self-Reinforcement}
\label{subsec:chain}

When item $\mathcal{I}_1$ comes into contact with $\mathcal{I}_2$, and
$\mathcal{I}_2$ is already in contact with $\mathcal{I}_3$, the following
occurs. The contact $\mathcal{I}_1$--$\mathcal{I}_2$ establishes a new local
negation relationship: $\mathcal{I}_1$ is now *not* $\mathcal{I}_2$
unambiguously. But this new relationship also sharpens the existing boundary
between $\mathcal{I}_2$ and $\mathcal{I}_3$: $\mathcal{I}_2$ is now not only
*not* $\mathcal{I}_3$ but *not* $\mathcal{I}_3$-while-in-contact-with-$\mathcal{I}_1$.
The contact chain acquires additional specificity at every node.

\begin{theorem}[Chain Self-Reinforcement]
\label{thm:chain_reinforce}
For a contact chain $\mathcal{I}_1 - \mathcal{I}_2 - \cdots - \mathcal{I}_n$,
adding a new contact $\mathcal{I}_0 - \mathcal{I}_1$ increases the definitional
precision of every item in the chain. Specifically, the partition depth of the
chain from $\mathcal{I}_0$ to $\mathcal{I}_n$ satisfies:
\begin{equation}
    D(\mathcal{I}_0, \mathcal{I}_n) = n
    \label{eq:chain_depth}
\end{equation}
and the addition of $\mathcal{I}_0$ adds one unit of definitional specificity
to each of $\mathcal{I}_1, \ldots, \mathcal{I}_n$ without adding any units of
global negation cost to any of them.
\end{theorem}

\begin{proof}
Before $\mathcal{I}_0$ is added, the chain $\mathcal{I}_1 - \cdots -
\mathcal{I}_n$ has each item defined by its local negation partners within the
chain plus its global negation against $\mathcal{U}$. After $\mathcal{I}_0$
contacts $\mathcal{I}_1$, the item $\mathcal{I}_1$ gains a new local negation
partner. The boundary $\beta(\mathcal{I}_1, \mathcal{I}_0) = \beta_{\min}$ is
established without changing any other boundary in the chain, since the
boundaries $\beta(\mathcal{I}_k, \mathcal{I}_{k+1})$ are already at $\beta_{\min}$
(by definition of contact). The only change is that the \emph{context} of each
item in the chain now includes one additional specific partner at $\mathcal{I}_1$.
This context propagates through the chain as causal residues, increasing
definitional precision at each node.
\end{proof}

\subsection{Normal Force as Floor Resistance}
\label{subsec:normal}

The contact state $\beta(\mathcal{I}_1, \mathcal{I}_2) = \beta_{\min}$ is
stable because any further reduction would require $\beta < \beta_{\min}$, which
is forbidden by Postulate~\ref{post:floor}. Attempting to drive the two items
below $\beta_{\min}$ requires introducing a restoring tendency --- the resistance
of the partition floor to being violated.

This restoring tendency is what classical mechanics calls the normal force. It
is not a separate force. It is the gradient of the floor:
\begin{equation}
    F_{\mathrm{normal}} = -\frac{\partial \beta_{\min}}{\partial D}
    \bigg|_{D \to 0^+}
    \label{eq:normal}
\end{equation}
The normal force is zero when $D > D_{\min}$ (the floor is not being approached)
and becomes unbounded as $D \to 0^+$ (the floor is impenetrable). This explains
why the normal force appears only at contact and not before: it is not a
pre-existing repulsion but the gradient of the floor, which is only encountered
when partition depth reaches its minimum.

\begin{remark}[Craters and Deformation]
\label{rem:craters}
When a stone strikes the ground at high velocity, the impact carries accumulated
partition count $\Delta M = \dot{M} \cdot v \cdot t$ in the direction of motion.
This partition count must be redistributed across the contact chain when the
floor is reached. The redistribution manifests as deformation: the ground's
partition structure is rearranged to accommodate the incoming partition count.
The crater is the residue of this rearrangement --- the new partition depth
minimum of the combined stone-ground system after the incoming $\Delta M$ has
been absorbed. The fact that craters form \emph{after} impact, rather than the
ground pre-emptively resisting, follows immediately: the floor is encountered at
the moment of contact, not before.
\end{remark}

% =============================================================================
% SECTION 3: EMERGENCE OF PHYSICAL QUANTITIES
% =============================================================================

\section{Emergence of Time, Mass, Charge, and Space}
\label{sec:emergence}

\subsection{The Time-Count Identity}
\label{subsec:time_count}

The partition count $M$ is defined as the cumulative number of completed
partition events in the history of a system. By Theorem~\ref{thm:floor}, $M$
is strictly monotone increasing. We now show that $M$ is proportional to
physical time.

\begin{theorem}[Time-Count Identity]
\label{thm:time_count}
For a system undergoing periodic partitioning with angular frequency $\omega$,
the rate of partition count accumulation is:
\begin{equation}
    \frac{dM}{dt} = \frac{\omega}{2\pi}
    \label{eq:time_count}
\end{equation}
Equivalently, the elapsed physical time between partition counts $M_1$ and
$M_2$ is:
\begin{equation}
    \Delta t = \frac{2\pi(M_2 - M_1)}{\omega} = \frac{\Delta M}{\dot{M}}
    \label{eq:time_from_count}
\end{equation}
\end{theorem}

\begin{proof}
Each partition event corresponds to one complete oscillatory cycle. One cycle
at frequency $\omega/(2\pi)$ Hz takes time $\tau_p = 2\pi/\omega$. Therefore
$\Delta M$ partition events take time $\Delta t = \Delta M \cdot \tau_p =
2\pi \Delta M / \omega$. Inverting gives $dM/dt = \omega/(2\pi)$.
\end{proof}

\begin{corollary}[Strict Monotonicity of Time]
\label{cor:time_mono}
Physical time, defined as $t = 2\pi M / \omega$, is strictly monotone
increasing. This follows directly from Theorem~\ref{thm:floor}: since $M$
cannot decrease, $t$ cannot decrease.
\end{corollary}

Corollary~\ref{cor:time_mono} is not a postulate about the arrow of time. It is
a theorem about the structure of partition events. The arrow of time does not
require an appeal to entropy, coarse-graining, or boundary conditions at the Big
Bang. It is a direct consequence of the irreducibility of the partition residue.

\subsection{Mass as Partition Inertia}
\label{subsec:mass}

\begin{definition}[Partition Inertia]
\label{def:inertia}
The \emph{partition inertia} of an item $\mathcal{I}$ is the quantity
$\mu(\mathcal{I})$ defined by:
\begin{equation}
    \mu = \alpha \cdot (m/z)
    \label{eq:inertia}
\end{equation}
where $m/z$ is the mass-to-charge ratio of $\mathcal{I}$, $\alpha$ is the
partition coupling constant, and the product $\mu$ determines how strongly the
item resists changes to its partition state.
\end{definition}

The physical mass $m$ of an item is its partition inertia in the limit of unit
charge ($z = 1$): $m = \mu/\alpha$. Mass is therefore not an intrinsic property
of matter but a measure of the item's resistance to partition state change ---
the cost, in partition count units, of transitioning from one partition state
$(n,\ell,m,s)$ to another.

\begin{theorem}[Mass-Frequency Relation]
\label{thm:mass_freq}
For an item with partition inertia $\mu$ undergoing oscillatory partition with
characteristic frequency $\omega_0$:
\begin{equation}
    m = \frac{\hbar \omega_0}{c^2}
    \label{eq:mass_freq}
\end{equation}
where $\hbar = E_{\mathrm{tick}}/(2\pi)$ is the reduced partition quantum and
$c = 2\pi$ rad/tick is the partition propagation rate.
\end{theorem}

\begin{proof}
The energy of one partition cycle is $E_{\mathrm{tick}} = \hbar \omega_0$.
By mass-energy equivalence (derived from the partition propagation rate $c$):
$E = mc^2$. Therefore $m = E/c^2 = \hbar\omega_0/c^2$.
\end{proof}

\subsection{Charge as Partition Asymmetry}
\label{subsec:charge}

The orientation index $m \in \{-\ell, \ldots, +\ell\}$ in the partition state
$(n,\ell,m,s)$ encodes the asymmetry of the boundary between $\mathcal{I}$ and
$\overline{\mathcal{I}}$. When $m = 0$, the boundary is symmetric: the item is
equally bounded on all sides. When $m \neq 0$, the boundary is asymmetric: the
item protrudes preferentially in one direction of the negation space.

\begin{definition}[Charge as Partition Orientation]
\label{def:charge}
The \emph{effective charge} of an item $\mathcal{I}$ in partition state
$(n,\ell,m,s)$ is:
\begin{equation}
    q = q_0 \cdot m
    \label{eq:charge_def}
\end{equation}
where $q_0$ is the elementary charge unit. Charge is the orientation of the
partition boundary --- its asymmetry relative to the global complement.
\end{definition}

\begin{theorem}[Charge Conservation]
\label{thm:charge_conserve}
For an isolated system of items $\{\mathcal{I}_k\}$ with total orientation
$\sum_k m_k = M_{\mathrm{orient}}$, the total charge is conserved:
\begin{equation}
    \sum_k q_k = q_0 M_{\mathrm{orient}} = \mathrm{const}
    \label{eq:charge_conserve}
\end{equation}
\end{theorem}

\begin{proof}
The orientation indices $\{m_k\}$ of the items in the system sum to
$M_{\mathrm{orient}}$, which is the net asymmetry of the system's boundary
with $\overline{\cup_k \mathcal{I}_k}$. Since the global complement is fixed
for an isolated system, and since the partition boundary between the system and
its complement has fixed total asymmetry, $M_{\mathrm{orient}}$ is constant.
\end{proof}

\begin{theorem}[Electromagnetic Interaction from Partition Orientation]
\label{thm:em_from_partition}
Two items $\mathcal{I}_1, \mathcal{I}_2$ with orientation indices $m_1, m_2$
interact through a potential:
\begin{equation}
    V(D) = \frac{q_0^2 m_1 m_2}{4\pi\epsilon_0} \cdot \frac{1}{r(D)}
    \label{eq:em_potential}
\end{equation}
where $r(D)$ is the spatial distance corresponding to partition depth $D$
(derived in Section~\ref{subsec:space}). Items with the same sign of $m$ have
parallel boundary orientations and repel (their local negation is less efficient
when they are close, since their boundaries point in the same direction and
compete). Items with opposite sign of $m$ have anti-parallel orientations and
attract (their local negation is more efficient when close, since each provides
what the other is negating).
\end{theorem}

\subsection{The S-Entropy Coordinates}
\label{subsec:sentropy}

For a partition state $(n,\ell,m,s)$, define the following normalised coordinates:

\begin{align}
    S_k &= \frac{n-1}{n_P - 1} \in [0,1]
    \label{eq:sk} \\
    S_t &= \frac{\ell}{n-1} \in [0,1] \quad (n > 1)
    \label{eq:st} \\
    S_e &= \frac{m + \ell}{2\ell} \in [0,1] \quad (\ell > 0)
    \label{eq:se}
\end{align}

where $n_P$ is the Planck depth (defined in Section~\ref{sec:composition}).
These three coordinates map every partition state to a point in $[0,1]^3$.

\begin{theorem}[S-Entropy as Partition Distance]
\label{thm:sentropy}
The S-entropy vector $\mathbf{S} = (S_k, S_t, S_e)$ measures the distance of
the partition state $(n,\ell,m,s)$ from the partition ground state $(1,0,0,s)$.
The Euclidean distance in $[0,1]^3$ is a well-defined metric on partition states.
\end{theorem}

\begin{proof}
The ground state $(1,0,0,s)$ has $S_k = 0$, $S_t = 0$, $S_e = 1/2$ (by
convention for the undefined case $\ell = 0$). Any other state has $S_k > 0$
and generally $S_t, S_e \neq 0, 1/2$. The Euclidean distance
$d = \sqrt{(S_k - 0)^2 + (S_t - 0)^2 + (S_e - 1/2)^2}$ is non-negative,
symmetric, and satisfies the triangle inequality, making it a metric.
\end{proof}

\subsection{Space as Emergent Partition Depth}
\label{subsec:space}

Space is not the container in which partition events occur. Space is the
coordinate representation of partition depth. We make this precise as follows.

\begin{definition}[Spatial Distance from Partition Depth]
\label{def:space}
The spatial distance $r(\mathcal{I}_1, \mathcal{I}_2)$ between two items is:
\begin{equation}
    r(\mathcal{I}_1, \mathcal{I}_2) = \lambda_P \cdot D(\mathcal{I}_1,
    \mathcal{I}_2)
    \label{eq:space}
\end{equation}
where $\lambda_P$ is the Planck length --- the spatial scale corresponding to
one unit of partition depth.
\end{definition}

This definition immediately gives:
\begin{enumerate}
    \item \textbf{Finite speed of signal propagation}: A signal propagating
          from $\mathcal{I}_1$ to $\mathcal{I}_2$ must traverse $D$ partition
          boundaries, each requiring time $\tau_p = 2\pi/\omega_0$ per boundary.
          Maximum speed: $c = \lambda_P/\tau_P = \lambda_P \omega_0/(2\pi)$.

    \item \textbf{Spatial dimensionality}: The three coordinates $(S_k, S_t,
          S_e)$ of S-entropy space define three independent directions of
          partition depth variation, giving rise to three spatial dimensions.

    \item \textbf{Curvature from partition density}: Regions with high
          partition count density (many items in close contact) have shorter
          effective $\lambda_P$ (more partition events per unit path length),
          manifesting as spacetime curvature in the coordinate representation.
\end{enumerate}

\subsection{Planck Depth and the Composition Limit}
\label{subsec:planck_depth}

\begin{definition}[Planck Depth]
\label{def:planck_depth}
For a system with $d$ distinguishable partition directions and oscillation
period $\tau_{\mathrm{osc}}$, the \emph{Planck depth} is:
\begin{equation}
    n_P = 1 + \left\lceil \log_{d+1} \left( \frac{\tau_{\mathrm{osc}}}{d \cdot
    \tau_P} \right) \right\rceil
    \label{eq:planck_depth}
\end{equation}
where $\tau_P = \sqrt{\hbar G / c^5}$ is the Planck time.
\end{definition}

For Cs-133 ($\tau_{\mathrm{osc}} = 1/9{,}192{,}631{,}770$ s, $d = 3$): $n_P = 56$.
This means the Cs-133 hyperfine transition operates at partition depth 56 --- it
requires 56 nested boundary layers to be well-defined as a distinct oscillatory
process. The partition count accumulated per second of Cs-133 oscillation is:
\begin{equation}
    \dot{M}_{\mathrm{Cs}} = 9{,}192{,}631{,}770 \;\mathrm{Hz}
    \label{eq:cs_mdot}
\end{equation}
which is by definition of the SI second.

\subsection{The Partition Lagrangian}
\label{subsec:lagrangian}

The dynamics of an item with partition inertia $\mu$ in a partition field are
governed by the Partition Lagrangian:

\begin{equation}
    \mathcal{L}_M = \frac{1}{2}\mu|\dot{x}|^2 + \mu\dot{x} \cdot A_M(x,t)
    - M(x,t)
    \label{eq:lagrangian}
\end{equation}

where $A_M(x,t)$ is the partition vector potential and $M(x,t)$ is the
partition count field. The Euler-Lagrange equations of
\eqref{eq:lagrangian} yield the equations of motion for the item in any
analyser geometry. We derive the four standard mass analyser equations as
consequences.

\begin{theorem}[Four Analyser Equations from Partition Lagrangian]
\label{thm:four_analysers}
The Euler-Lagrange equations of $\mathcal{L}_M$ in the appropriate boundary
conditions yield:
\begin{align}
    \text{TOF:} \quad T_{\mathrm{TOF}} &= L\sqrt{\frac{m}{2qV}}
    \label{eq:tof} \\
    \text{Orbitrap:} \quad \omega_z &= \sqrt{\frac{q\kappa}{m}}
    \label{eq:orbitrap} \\
    \text{FT-ICR:} \quad \omega_c &= \frac{qB}{m}
    \label{eq:fticr} \\
    \text{Quadrupole:} \quad a_z &= -2q_z \omega^2 / 4
    \label{eq:quad}
\end{align}
where $L$ is the TOF flight path length, $V$ is the accelerating voltage,
$\kappa$ is the Orbitrap geometry constant, $B$ is the FT-ICR magnetic field,
and $q_z$ is the Mathieu stability parameter.
\end{theorem}

\begin{proof}[Proof sketch]
Each analyser imposes a specific boundary condition on the partition field
$M(x,t)$. The TOF imposes a linear potential $M = qVx/L$; the Orbitrap imposes
a quadratic potential $M = \kappa x^2/2$; the FT-ICR imposes a vector potential
$A_M = B \times x/2$; the quadrupole imposes an oscillating potential
$M = (a - 2q\cos(2\Omega t))x^2/4$. Substituting each into the Euler-Lagrange
equations of \eqref{eq:lagrangian} and solving gives the respective frequency
or transit time equations above. Full derivations are given in
Appendix~\ref{app:analyser_derivations}.
\end{proof}

\subsection{The Partition Rates of the Four Analysers}
\label{subsec:partition_rates}

From Theorem~\ref{thm:time_count} and Theorem~\ref{thm:four_analysers}, the
partition rate $\dot{M} = \omega/(2\pi)$ for each analyser is:

\begin{align}
    \dot{M}_{\mathrm{Orb}}(m/z) &= \frac{1}{2\pi}\sqrt{\frac{q\kappa}{m}}
    \label{eq:mdot_orb} \\
    \dot{M}_{\mathrm{ICR}}(m/z) &= \frac{qB}{2\pi m}
    \label{eq:mdot_icr} \\
    \dot{M}_{\mathrm{TOF}}(m/z) &= \frac{1}{L}\sqrt{\frac{2qV}{m}} \cdot
    \frac{1}{2\pi}
    \label{eq:mdot_tof} \\
    \dot{M}_{\mathrm{quad}} &= \frac{q_u \Omega}{4\pi}
    \label{eq:mdot_quad}
\end{align}

Note that $\dot{M}_{\mathrm{quad}}$ is mass-independent at the stability apex
($q_u$ fixed), while all others scale as $(m/z)^{-1/2}$ (Orbitrap, TOF) or
$(m/z)^{-1}$ (FT-ICR). These scaling laws are direct predictions of the
Partition Lagrangian with zero free parameters.

% =============================================================================
% SECTION 4: COMPOSITION INFLATION AND THE RESIDUE AS CAUSAL PROPAGATOR
% =============================================================================

\section{Composition Inflation, the Causal Residue, and the Origin of the
Second}
\label{sec:composition}

\subsection{Composition Inflation}
\label{subsec:comp_inflation}

Each partition event is not just a state transition --- it is a branching of the
space of possible next states. The partition count $M$ advances by 1, but the
number of distinguishable ways it can advance depends on the number of available
partition directions $d$.

\begin{theorem}[Composition Inflation]
\label{thm:comp_inflation}
For a system with $d$ distinguishable partition directions, the number of
distinguishable categorical trajectories of length $n$ is:
\begin{equation}
    T(n,d) = d(d+1)^{n-1}
    \label{eq:comp_inflation}
\end{equation}
\end{theorem}

\begin{proof}
At depth 1, there are $d$ distinguishable starting directions: $T(1,d) = d$.
At each subsequent depth, each existing trajectory can branch into $d+1$
directions (the $d$ original directions plus the ``stay'' option, which
corresponds to remaining in the current partition state). Therefore
$T(n,d) = d(d+1)^{n-1}$.
\end{proof}

For $d=3$ (three spatial partition directions): $T(n,3) = 3 \cdot 4^{n-1}$.
For Cs-133 at $n_P = 56$: $T(56,3) = 3 \cdot 4^{55} = 3.81 \times 10^{33}$,
exceeding the $2.02 \times 10^{33}$ Planck intervals per Cs-133 oscillation
period. The partition framework generates more distinguishable trajectories than
there are Planck intervals, confirming that the Cs-133 oscillation is a
legitimate classical process fully resolved by the partition structure at
depth 56.

\subsection{The Residue as Causal Propagator}
\label{subsec:causal_residue}

We now establish the central claim of this paper.

\begin{theorem}[Residue as Causal Propagator]
\label{thm:causal_residue}
The irreducible residue $\beta > 0$ of partition $k$ is the necessary and
sufficient condition for the existence of partition $k+1$. Without the residue,
there is no causal link between successive partitions, and the universe has no
mechanism to progress from one state to the next.
\end{theorem}

\begin{proof}
Partition $k$ establishes the distinction $\mathcal{I} \neq \overline{\mathcal{I}}$
and generates residue $\beta_k > 0$. This residue is deposited at the boundary
between $\mathcal{I}$ and $\overline{\mathcal{I}}$. The residue constitutes a
physical difference between the state of the system immediately before partition
$k$ and immediately after. Partition $k+1$ occurs at the boundary, driven by the
presence of $\beta_k$: the residue creates a non-equilibrium condition at the
boundary, which is the driving condition for the next partition event. If
$\beta_k = 0$, no physical difference would exist between the pre- and
post-partition states, no non-equilibrium condition would exist, and partition
$k+1$ would have no cause. The residue is therefore the causal link. It is
necessary (no residue implies no next partition) and sufficient (any non-zero
residue creates the driving condition for the next partition).
\end{proof}

\begin{corollary}[The Three Forms of the Residue]
\label{cor:three_forms}
The residue $\beta_k$ manifests in three forms depending on which aspect of
the next partition is being described:
\begin{enumerate}
    \item \textbf{As time}: $\beta_k$ determines the interval $\tau_{p,k+1}$
          of the next partition event. Time is accumulated residue of partitioning.

    \item \textbf{As entropy}: When many residues $\{\beta_k\}$ are present and
          their specific causal links to future partitions are not individually
          tracked, the aggregate is entropy. Shannon entropy
          $H = -\sum_k p_k \log p_k$ is the information content of the
          residue distribution when individual causal links are unresolved.

    \item \textbf{As information}: When the residue $\beta_k$ from partition $k$
          is the specific input to partition $k+1$ of a different item, the
          residue is \emph{information} --- the structural imprint of one
          partition on another's boundary conditions.
\end{enumerate}
Time, entropy, and information are not three different things. They are the
same residue described from three different perspectives on the next partition.
\end{corollary}

\subsection{The Angular Reformulation of Fundamental Constants}
\label{subsec:angular}

The composition inflation framework reveals the partition-theoretic origin of
fundamental constants. We state the key reductions:

\begin{align}
    c &= 2\pi \;\mathrm{rad/tick}
    \label{eq:c_angular} \\
    \hbar &= \frac{E_{\mathrm{tick}}}{2\pi}
    \label{eq:hbar_angular} \\
    k_B &= \frac{E_{\mathrm{tick}}}{\ln(d+1)}
    \label{eq:kb_angular}
\end{align}

where $E_{\mathrm{tick}}$ is the energy per partition tick (the minimum quantum
of action in the partition framework) and $d+1 = 4$ for three spatial dimensions.

\begin{theorem}[G is Irreducible]
\label{thm:G_irreducible}
The gravitational constant $G$ cannot be expressed in terms of $E_{\mathrm{tick}}$,
$d$, $\tau_P$, or any combination of partition counting quantities.
\end{theorem}

\begin{proof}
Every reducible constant ($c$, $\hbar$, $k_B$) describes the rate or quantum of
a partition process: $c$ is a propagation rate, $\hbar$ is an energy per radian,
$k_B$ is an energy per log-state. All are expressible as $E_{\mathrm{tick}}$
times a dimensionless counting factor. $G$ appears in the gravitational
acceleration $a = GM/r^2$, where $M$ is partition inertia (mass) and $r$ is
spatial distance (partition depth times $\lambda_P$). Substituting:
\begin{equation}
    a = G \cdot \frac{\mu/\alpha}{\lambda_P^2 \cdot D^2}
\end{equation}
For $a$ to be expressed in terms of partition counting quantities, $G$ would
need to reduce to $[E_{\mathrm{tick}}] \cdot [\mathrm{dimensionless}]$ divided
by the units of $\mu/(\alpha \lambda_P^2)$. But $\lambda_P$ is itself defined in
terms of $G$: $\lambda_P = \sqrt{\hbar G/c^3}$. Substituting this definition
into the expression for $a$ makes $G$ appear on both sides of the equation in a
form that does not cancel. No rearrangement eliminates $G$. It is therefore
irreducible.
\end{proof}

The irreducibility of $G$ reflects the partition-theoretic insight of
Remark~\ref{rem:G_irreducible}: $G$ is the contact geometry constant, not a
process constant. Contact --- the ground state of partition depth minimisation
--- is prior to the processes that approach it, and therefore $G$ is prior to
the constants that describe those processes.

\subsection{The Partition Second}
\label{subsec:partition_second}

\begin{definition}[Partition Second]
\label{def:partition_second}
The \emph{partition second} (Ps) is the elapsed duration expressed as a ratio
of partition counts:
\begin{equation}
    \Delta t \;[\mathrm{Ps}] = \frac{\Delta M_{\mathcal{A}}}{\dot{M}_{\mathcal{A}}
    \cdot \dot{M}_{\mathrm{Cs}}} = \frac{\Delta t \;[\mathrm{s}]}{1\;\mathrm{s}}
    \label{eq:partition_second}
\end{equation}
where $\mathcal{A}$ is any analyser and $\dot{M}_{\mathrm{Cs}} = 9{,}192{,}631{,}770$
Hz is the Cs-133 hyperfine rate.
\end{definition}

\begin{theorem}[Universal Conversion Constant]
\label{thm:K}
For an Orbitrap with geometry constant $\kappa$, the number of Cs-133 hyperfine
cycles per Orbitrap axial cycle at mass-to-charge ratio $m/z$ is:
\begin{equation}
    \frac{\dot{M}_{\mathrm{Cs}}}{\dot{M}_{\mathrm{Orb}}(m/z)} = K \cdot
    \sqrt{m/z}
    \label{eq:K_conversion}
\end{equation}
where:
\begin{equation}
    K = \frac{\dot{M}_{\mathrm{Cs}} \cdot 2\pi}{\sqrt{e\kappa/u}}
    = 588.016 \;\mathrm{Cs\;cycles\;Orbitrap\;cycle}^{-1}\;u^{-1/2}
    \label{eq:K_value}
\end{equation}
is a universal constant derivable from CODATA values alone.
\end{theorem}

\begin{proof}
From Eq.~\eqref{eq:mdot_orb}: $\dot{M}_{\mathrm{Orb}} = \sqrt{q\kappa/m}/(2\pi)$.
Therefore:
\begin{equation}
    \frac{\dot{M}_{\mathrm{Cs}}}{\dot{M}_{\mathrm{Orb}}} =
    \frac{\dot{M}_{\mathrm{Cs}} \cdot 2\pi}{\sqrt{q\kappa/m}} =
    \frac{\dot{M}_{\mathrm{Cs}} \cdot 2\pi}{\sqrt{e\kappa/u}} \cdot \sqrt{m/z}
    = K\sqrt{m/z}
\end{equation}
where we have used $q = e \cdot z$ and $m = u \cdot (m/z)$ with $e$ the
elementary charge and $u$ the atomic mass unit.
\end{proof}

\begin{corollary}[Analyser-Independent Time]
\label{cor:analyser_independent}
For any analyser $\mathcal{A}$ with partition rate $\dot{M}_{\mathcal{A}}$:
\begin{equation}
    \Delta t \;[\mathrm{s}] = \frac{\Delta M_{\mathcal{A}}}{\dot{M}_{\mathcal{A}}}
    \label{eq:analyser_independent}
\end{equation}
holds independently of which analyser is used. Two analysers with different
$\dot{M}$ accumulate different $\Delta M$ over the same duration, but the ratio
$\Delta M_{\mathcal{A}}/\dot{M}_{\mathcal{A}}$ is identical. The partition
second is analyser-invariant.
\end{corollary}

\subsection{The Analyser-Switching Time Jump}
\label{subsec:time_jump}

\begin{definition}[Partition Rate Discontinuity]
\label{def:rate_discontinuity}
When an ion transitions from analyser $\mathcal{A}_1$ to analyser $\mathcal{A}_2$
at physical time $t_s$, the partition rate changes discontinuously:
\begin{equation}
    \dot{M}(t) = \begin{cases}
        \dot{M}_{\mathcal{A}_1} & t < t_s \\
        \dot{M}_{\mathcal{A}_2} & t > t_s
    \end{cases}
    \label{eq:rate_discontinuity}
\end{equation}
The partition count $M(t)$ is continuous at $t_s$ but its rate is not.
\end{definition}

\begin{definition}[Time Jump]
\label{def:time_jump}
The \emph{time jump} $\Delta M_{\mathrm{jump}}$ is the difference in partition
count accumulated over interval $[t_s, t_s + \tau]$ between the two analysers:
\begin{equation}
    \Delta M_{\mathrm{jump}}(\tau) = \left(\dot{M}_{\mathcal{A}_2} -
    \dot{M}_{\mathcal{A}_1}\right) \cdot \tau
    \label{eq:time_jump}
\end{equation}
The time jump is the measurable signature of the metric discontinuity at the
analyser boundary --- the residue of the partition between the two analyser
regimes.
\end{definition}

\begin{theorem}[Time Jump Scaling]
\label{thm:jump_scaling}
For the Q-Orbitrap DDA cycle (quadrupole MS1, Orbitrap HCD), the time jump
$\Delta M_{\mathrm{jump}}$ for an ion at $m/z$ over MS1 window $t_{\mathrm{MS1}}$
is:
\begin{equation}
    \Delta M_{\mathrm{jump}}(m/z) = \left(\frac{\sqrt{q\kappa/m}}{2\pi} -
    \frac{q_u\Omega}{4\pi}\right) \cdot t_{\mathrm{MS1}}
    \label{eq:jump_scaling}
\end{equation}
This quantity scales as $(m/z)^{-1/2}$ for the Orbitrap term and is constant
for the quadrupole term. Therefore $\Delta M_{\mathrm{jump}}$ is a
monotone-decreasing function of $m/z$, larger for light ions and smaller for
heavy ions.
\end{theorem}

\subsection{Partition Gradient Trajectories}
\label{subsec:gradient_traj}

\begin{definition}[Partition Gradient Trajectory]
\label{def:gradient_traj}
For an ion at $m/z$ observed at retention times $\{t_k\}$ along a
chromatographic elution, the \emph{partition gradient trajectory} is the sequence:
\begin{equation}
    \{M_{\mathcal{A}}(t_k)\} = \{\dot{M}_{\mathcal{A}}(m/z) \cdot t_k\}
    \label{eq:gradient_traj}
\end{equation}
for analyser $\mathcal{A}$. Different analysers produce different trajectories
for the same ion at the same retention times, with slopes $\dot{M}_{\mathcal{A}}(m/z)$.
\end{definition}

\begin{theorem}[Retention Time from Partition Counts]
\label{thm:rt_from_partition}
For two observations of the same ion at $m/z$ at retention times $t_i$ and
$t_j$, the ratio of retention times equals the ratio of accumulated partition
counts:
\begin{equation}
    \frac{t_i}{t_j} = \frac{M_{\mathcal{A}}(t_i)}{M_{\mathcal{A}}(t_j)}
    \label{eq:rt_ratio}
\end{equation}
This holds for any analyser $\mathcal{A}$, since the same $\dot{M}_{\mathcal{A}}$
cancels from numerator and denominator.
\end{theorem}

\begin{corollary}[DDA Cycle Count from Partition Accumulation]
\label{cor:dda_cycles}
The number of complete DDA cycles elapsed before retention time $t$ is:
\begin{equation}
    n_{\mathrm{cycles}}(t) = \frac{t}{t_{\mathrm{cycle}}} =
    \frac{M_{\mathrm{total}}(t)}{M_{\mathrm{cycle}}}
    \label{eq:dda_cycles}
\end{equation}
where $M_{\mathrm{total}}(t) = \bar{\dot{M}} \cdot t$ is the total partition
count at time $t$ using the time-averaged partition rate $\bar{\dot{M}} =
\dot{M}_{\mathrm{quad}} \cdot f_{\mathrm{quad}} + \dot{M}_{\mathrm{Orb}} \cdot
f_{\mathrm{Orb}}$, and $M_{\mathrm{cycle}} = \bar{\dot{M}} \cdot t_{\mathrm{cycle}}$.
The number of analyser-switching events is therefore directly readable from the
partition accumulation without reference to any clock.
\end{corollary}

% =============================================================================
% SECTION 5: THE LOSCHMIDT PARADOX AS CATEGORY ERROR
% =============================================================================

\section{The Loschmidt Paradox as a Category Error}
\label{sec:loschmidt}

\subsection{Statement of the Paradox}
\label{subsec:paradox_statement}

Loschmidt's paradox arises from the following argument. The microscopic equations
of motion of classical and quantum mechanics are time-reversal symmetric: if
$\{x_i(t), p_i(t)\}$ is a solution, then $\{x_i(-t), -p_i(-t)\}$ is also a
solution. Therefore, if a system evolves from state $A$ to state $B$, it should
be possible to reverse all momenta at state $B$ and recover state $A$. This
appears to contradict the second law of thermodynamics, which asserts that
macroscopic processes are irreversible.

Conventional resolutions appeal to coarse-graining, the overwhelming probability
of high-entropy states, or special initial conditions at the Big Bang. We argue
that all such resolutions share a common flaw: they accept the premise of the
paradox --- that the velocity reversal operation is coherently definable --- and
then argue about probabilities. We instead show that the premise is incoherent,
and that the paradox dissolves before any probabilistic argument is needed.

\subsection{The Frozen Universe Requirement}
\label{subsec:frozen}

\begin{theorem}[Velocity Reversal Requires a Frozen Universe]
\label{thm:frozen}
The Loschmidt velocity reversal operation is coherently definable only in a
universe where the complement of the system under study is frozen --- that is,
does not progress --- during the reversal operation. A frozen universe is
impossible under the partition framework.
\end{theorem}

\begin{proof}
The velocity reversal operation consists of:
\begin{enumerate}
    \item Identifying the system $\mathcal{S}$ and its state $\{x_i(t_s),
          p_i(t_s)\}$ at reversal time $t_s$
    \item Constructing the reversed state $\{x_i(t_s), -p_i(t_s)\}$
    \item Releasing the system into this reversed state
\end{enumerate}
Step 1 requires observing the system, which is a partition event: $\mathcal{S}$
must be distinguished from $\overline{\mathcal{S}}$ to be observed. This
partition event takes time $\tau_{p,\mathrm{obs}} > 0$
(Postulate~\ref{post:duration}) and generates residue
$\beta_{\mathrm{obs}} > 0$ (Postulate~\ref{post:floor}).

During $\tau_{p,\mathrm{obs}}$, the complement $\overline{\mathcal{S}}$ must
progress (Postulate~\ref{post:duration}), advancing from state $U(t_s)$ to
$U(t_s + \tau_{p,\mathrm{obs}})$.

Step 2 requires computing the negated momenta, which takes additional time
$\tau_{p,\mathrm{compute}} > 0$ during which $\overline{\mathcal{S}}$ advances
to $U(t_s + \tau_{p,\mathrm{obs}} + \tau_{p,\mathrm{compute}})$.

Step 3 releases $\mathcal{S}$ into state $\{x_i(t_s), -p_i(t_s)\}$ --- a state
computed at $t_s$ --- into a universe now at
$U(t_s + \tau_{p,\mathrm{obs}} + \tau_{p,\mathrm{compute}})$, not $U(t_s)$. The
system is released into the wrong universe. The reversed state is only
time-reversed relative to $U(t_s)$; the universe has moved on to
$U(t_s + \Delta\tau)$ with $\Delta\tau > 0$.

For the reversal to be coherent, we would need $\overline{\mathcal{S}}$ to
remain at $U(t_s)$ throughout. But a frozen $\overline{\mathcal{S}}$ violates
Postulate~\ref{post:duration}: if $\overline{\mathcal{S}}$ does not progress,
no partition event within it can be well-defined, and the observation in Step 1
cannot be performed. The frozen universe is therefore impossible under the
partition framework, and the velocity reversal operation is incoherent.
\end{proof}

\subsection{The Swimmer Argument}
\label{subsec:swimmer}

We illustrate the abstract argument with a concrete case. Consider a swimmer
performing the backstroke in a pool. The system $\mathcal{S}$ is the
swimmer-pool system. The Loschmidt argument claims: reverse all velocities, and
the swimmer will reverse their trajectory back to the starting position.

\begin{theorem}[Causal Path Removal by Reversal]
\label{thm:swimmer}
Reversing all interactions in the swimmer-pool system does not time-reverse the
system; it removes the causal path that constituted the system's history.
\end{theorem}

\begin{proof}
The swimmer's stroke at time $t_0$ generates a pressure wave in the pool. This
pressure wave is a residue $\beta_{\mathrm{stroke}}$ deposited into the
pool-water contact chain. The pressure wave:
\begin{enumerate}
    \item Propagates to the pool walls at time $t_0 + L/c_{\mathrm{water}}$,
          generating residue $\beta_{\mathrm{wall}}$
    \item Reflects and propagates to the far side, generating residue
          $\beta_{\mathrm{reflect}}$
    \item Deposits thermal energy throughout the pool, generating residues
          $\{\beta_{\mathrm{thermal},k}\}$
    \item Propagates into the pool structure, the ground, and the
          surrounding environment, generating further residues
\end{enumerate}
Each residue $\beta_k$ is the causal propagator for the next event in the chain
(Theorem~\ref{thm:causal_residue}). The causal path of the stroke is the
sequence $\beta_{\mathrm{stroke}} \to \beta_{\mathrm{wall}} \to
\beta_{\mathrm{reflect}} \to \cdots$.

By time $t_1 > t_0$, the residues $\{\beta_k\}$ have propagated throughout the
contact chain of the universe. They are no longer localised in the pool. They
are in the pool structure, the ground, the air, the building, the Earth's crust.

To reverse the swimmer, one would need to reverse all momenta including those
of all items that have received any $\beta_k$. This includes the pool walls,
the pool structure, the ground, the air, and the Earth's crust. The reversal
requirement expands without bound through the contact chain.

Attempting to reverse a finite subset --- the pool and swimmer but not the
ground --- would be incomplete: the ground's contact chain now includes the
residues of the stroke and has advanced its partition count $M_{\mathrm{ground}}$
accordingly. The pool, once reversed, would be embedded in a ground that has
already received and propagated the stroke's residues. The contact chain between
pool and ground would generate forward-directed partition events immediately upon
contact, because the ground's state is not time-reversed.

Most fundamentally: the causal path of the stroke --- the sequence of residues
$\beta_k$ that propagated causally from the stroke --- cannot be reversed because
reversing it would require each residue to un-cause the next one. A residue
cannot un-cause a subsequent partition event: the partition event has already
established new contact relationships that themselves have residues. Un-causation
is not a physical operation; it would require the partition count $M$ to
decrease, which Theorem~\ref{thm:floor} forbids.

Therefore reversing all velocities does not time-reverse the system. It removes
the causal path that made the system a system.
\end{proof}

\subsection{The Spin Echo: A Forward Process Misread as Reversal}
\label{subsec:spin_echo}

The nuclear magnetic resonance spin echo appears to support time-reversal: after
a $\pi$ pulse at time $\tau$, the spins refocus at time $2\tau$ as if the
dephasing had been reversed. We show this is a forward process.

\begin{theorem}[Spin Echo as Forward Partition Process]
\label{thm:spin_echo}
The NMR spin echo is not a time-reversed process. It is a forward partition
process whose initial conditions were set by the residue of the $\pi$ pulse
interacting with the residue of the dephasing process.
\end{theorem}

\begin{proof}
The spin echo proceeds as follows in partition terms:
\begin{enumerate}
    \item \textbf{Excitation}: $\pi/2$ pulse establishes phase coherence.
          Partition count $M_0$.
    \item \textbf{Dephasing} (interval $[0,\tau]$): Each spin $k$ accumulates
          partition count at rate $\dot{M}_k = \omega_k/(2\pi)$ where $\omega_k
          = \omega_0 + \delta\omega_k$ includes local field inhomogeneity
          $\delta\omega_k$. At time $\tau$: partition count of spin $k$ is
          $M_k(\tau) = \dot{M}_k \cdot \tau$. This is a forward process.
    \item \textbf{$\pi$ pulse at $\tau$}: Inverts $z$-component of each spin.
          In partition terms: applies a parity transformation to the residue
          parity $s \to -s$ of each spin's partition state. This is a new
          partition event with residue $\beta_\pi > 0$. Partition count advances
          to $M(\tau) + 1$.
    \item \textbf{Refocusing} (interval $[\tau, 2\tau]$): Each spin continues
          accumulating partition count forward at the same rate $\dot{M}_k$.
          At time $2\tau$: partition count is $M_k(2\tau) = \dot{M}_k \cdot
          2\tau$. The phase of spin $k$ is:
          \begin{equation}
              \phi_k(2\tau) = \omega_k \cdot 2\tau - 2\omega_k\tau = 0
          \end{equation}
          (the $\pi$ pulse effectively negated the accumulated phase by flipping
          the partition parity, and the subsequent forward accumulation cancels
          the dephasing exactly).
\end{enumerate}
The refocusing is not time reversal. The partition count $M_k$ has increased
monotonically from $0$ to $\dot{M}_k \cdot 2\tau$ throughout the entire echo
sequence. Time has advanced by $2\tau$, not returned to zero. The coherence
restoration arises because the $\pi$ pulse is a new partition event whose
residue (the parity flip) is precisely structured to cancel the accumulated
phase --- not to reverse it.

The universe advanced its partition count by $\dot{M}_{\mathrm{universe}} \cdot
2\tau$ throughout the echo. The magnet, the sample tube, the laboratory, the
Earth, all progressed. The echo occurred inside a universe that had moved
forward by $2\tau$, not inside the universe at time 0.
\end{proof}

\subsection{The T2* and T2 Distinction as Partition Depth Split}
\label{subsec:T2_split}

The distinction between $T_2^*$ (apparent transverse relaxation time, including
field inhomogeneity) and $T_2$ (true transverse relaxation time) maps directly
onto the partition framework:

\begin{itemize}
    \item $T_2^*$ is the dephasing time at the \emph{microstate} level: each
          spin's individual partition trajectory diverges due to local field
          variations. These are reversible in the sense that the $\pi$ pulse
          can cancel them --- they are phase errors, not entropy, because the
          local field variations are themselves part of the partition chain and
          their residues are coherent.

    \item $T_2$ is the dephasing time at the \emph{partition count} level:
          irreversible loss of coherence due to spin-spin interactions that
          deposit residues $\beta_k$ into the contact chain between spins.
          These residues propagate into the broader contact chain and cannot
          be recovered by the $\pi$ pulse.
\end{itemize}

The ratio $T_2^*/T_2$ is the ratio of reversible phase accumulation (microstate
level, recoverable by parity flip) to irreversible partition count accumulation
(contact chain level, not recoverable). The spin echo recovers $T_2^*$; it
cannot recover $T_2$.

\subsection{Formal Resolution of the Paradox}
\label{subsec:formal_resolution}

\begin{theorem}[Loschmidt Paradox Resolution]
\label{thm:resolution}
There is no paradox between the time-reversal symmetry of microscopic equations
and the irreversibility of macroscopic processes. The apparent paradox arises
from a category error: the time-reversal symmetry of equations of motion applies
to the \emph{form} of the equations, not to the \emph{states} to which those
equations apply. The states are constrained by the partition framework to satisfy
$M(t_2) > M(t_1)$ for $t_2 > t_1$. No state trajectory with decreasing $M$ is
physically accessible, regardless of whether the equations of motion are
time-reversal symmetric.
\end{theorem}

\begin{proof}
Let $\mathcal{E}$ be the equations of motion of a system $\mathcal{S}$, and let
$T$ be the time-reversal operator: $T: (x(t), p(t)) \mapsto (x(-t), -p(-t))$.
The statement ``$\mathcal{E}$ is time-reversal symmetric'' means: if
$(x(t), p(t))$ satisfies $\mathcal{E}$, then $(x(-t), -p(-t))$ also satisfies
$\mathcal{E}$.

This is a statement about the \emph{solution space} of $\mathcal{E}$: both
$(x(t), p(t))$ and its time-reverse are mathematical solutions. It is not a
statement that the time-reversed solution is physically accessible from the
forward solution.

Physical accessibility requires:
\begin{enumerate}
    \item The initial conditions of the reversed solution must be achievable by
          a physical operation. This requires the velocity reversal operation,
          which Theorem~\ref{thm:frozen} shows is incoherent.
    \item The reversed solution must have $M_{\mathrm{reversed}}(t_2) >
          M_{\mathrm{reversed}}(t_1)$ for $t_2 > t_1$. A time-reversed solution
          $(x(-t), -p(-t))$ has decreasing coordinate time, which corresponds
          to decreasing partition count. This violates Theorem~\ref{thm:floor}.
\end{enumerate}
Therefore, while time-reversed solutions exist as mathematical objects
satisfying $\mathcal{E}$, no physical process can access them. The paradox
arises from confusing mathematical solutions with physical states. There is no
paradox about physical states: all physical states have monotone increasing $M$,
and no physical operation can violate this.
\end{proof}

% =============================================================================
% SECTION 6: EXPERIMENTAL VALIDATION
% =============================================================================

\section{Experimental Validation}
\label{sec:experiment}

\subsection{Dataset and Protocol}
\label{subsec:dataset}

We validate the partition framework on the NIST Acylcarnitine and Carnitine
(AC/CAC) MS Library 2020, Version 1D1B \cite{nist_accac_2020}. This library
contains 5,328 MS/MS spectra of 244 unique precursor ions ($m/z$ range:
162.1125--650.3746 u), acquired by LC-ESI-MS/MS in positive ion mode with HCD
fragmentation on a Thermo Fisher Q-Orbitrap instrument. Instrument parameters:
quadrupole RF frequency $\Omega = 2\pi \times 1.1 \times 10^6$ rad s$^{-1}$,
stability apex $q_u = 0.706$, Orbitrap geometry constant $\kappa = 1.0 \times
10^8$ m$^{-2}$ \cite{makarov2000}. Typical DDA cycle: $t_{\mathrm{MS1}} = 50$
ms (quadrupole), $t_{\mathrm{HCD}} = 96$ ms (Orbitrap HCD transient).

The DDA cycle is a controlled, sequential partition process:
\begin{enumerate}
    \item \textbf{Quadrupole selection}: partitions the ion beam at
          $\dot{M}_{\mathrm{quad}}$, accumulating $\Delta M_{\mathrm{quad}} =
          \dot{M}_{\mathrm{quad}} \cdot t_{\mathrm{MS1}}$ partition counts
    \item \textbf{Orbitrap detection}: partitions the selected ions at
          $\dot{M}_{\mathrm{Orb}}(m/z)$, accumulating $\Delta M_{\mathrm{Orb}}
          = \dot{M}_{\mathrm{Orb}}(m/z) \cdot t_{\mathrm{HCD}}$ partition counts
    \item \textbf{Handoff}: the boundary between stages 1 and 2 carries a
          partition rate discontinuity $\dot{M}_{\mathrm{Orb}} -
          \dot{M}_{\mathrm{quad}}$, manifesting as a time jump $\Delta
          M_{\mathrm{jump}}$
\end{enumerate}

All quantities are computed analytically from the Partition Lagrangian with no
free parameters. Physical constants from CODATA 2018 \cite{codata2018}.

\subsection{Mass-Scaling Law: $(m/z)^{-1/2}$}
\label{subsec:scaling_law}

The Partition Lagrangian predicts $\dot{M}_{\mathrm{Orb}} \propto (m/z)^{-1/2}$
(Eq.~\eqref{eq:mdot_orb}) and $\dot{M}_{\mathrm{quad}} = \mathrm{const}$
(Eq.~\eqref{eq:mdot_quad}). Therefore the ratio $\dot{M}_{\mathrm{Orb}}/
\dot{M}_{\mathrm{quad}} \propto (m/z)^{-1/2}$, giving a log-log slope of
$-0.5$ exactly.

\textbf{Result}: Across all 244 unique $m/z$ values, the log-log slope of
$\dot{M}_{\mathrm{Orb}}(m/z)/\dot{M}_{\mathrm{quad}}$ vs $m/z$ is:
\begin{equation}
    \mathrm{slope}_{\mathrm{obs}} = -0.5000, \quad
    \mathrm{slope}_{\mathrm{pred}} = -0.5000, \quad
    |\mathrm{deviation}| < 10^{-15}
    \label{eq:slope_result}
\end{equation}
The deviation is at floating-point machine precision. The $m^{-1/2}$ law holds
exactly across the full mass range 162--650 u, confirming the Partition
Lagrangian's prediction for the Orbitrap axial frequency with zero free
parameters.

\subsection{Analyser-Switching Time Jumps}
\label{subsec:jump_results}

From Theorem~\ref{thm:jump_scaling}, the time jump for $t_{\mathrm{MS1}} = 50$
ms is:

\begin{equation}
    \Delta M_{\mathrm{jump}}(m/z) = \left(\dot{M}_{\mathrm{Orb}}(m/z) -
    \dot{M}_{\mathrm{quad}}\right) \times 0.050 \;\mathrm{s}
    \label{eq:jump_result}
\end{equation}

\textbf{Results} across all 5,328 spectra:
\begin{itemize}
    \item Range: $1.12 \times 10^4$ (at $m/z = 650.4$) to $4.20 \times 10^4$
          (at $m/z = 162.1$) partition counts
    \item As fraction of Orbitrap transient: mean $28.3\% \pm 2.1\%$
    \item Monotone decrease with $m/z$: the time jump is larger for lighter ions
          because $\dot{M}_{\mathrm{Orb}}$ is larger relative to
          $\dot{M}_{\mathrm{quad}}$ at low $m/z$
\end{itemize}

The time jump is not an instrumental artefact. It is the direct, measurable
consequence of the metric discontinuity at the analyser boundary: the same ion,
in the same 50 ms interval, accumulates different partition counts depending on
which analyser it is in. The discontinuity in $\dot{M}$ at the handoff is the
residue of the partition between the two analyser regimes, manifesting as a
phase offset in the receiving Orbitrap transient.

\subsection{Universal Conversion Constant K}
\label{subsec:K_result}

Theorem~\ref{thm:K} predicts:
\begin{equation}
    K = \frac{\dot{M}_{\mathrm{Cs}} \cdot 2\pi}{\sqrt{e\kappa/u}} =
    588.016 \;\mathrm{Cs\;cyc\;Orb\;cyc}^{-1}\;u^{-1/2}
    \label{eq:K_result}
\end{equation}

\textbf{Validation}: For each of the 109 unique $m/z$ values in the first 500
spectra, the observed ratio $\dot{M}_{\mathrm{Cs}}/\dot{M}_{\mathrm{Orb}}(m/z)
/ \sqrt{m/z}$ was computed. The deviation from $K = 588.016$ is:
\begin{equation}
    \max_{m/z} \left|\frac{\dot{M}_{\mathrm{Cs}}/\dot{M}_{\mathrm{Orb}}(m/z)
    / \sqrt{m/z} - K}{K}\right| \times 10^6 = 0.0000 \;\mathrm{ppm}
    \label{eq:K_deviation}
\end{equation}
The deviation is zero to machine precision across all 109 values. The Cs-133
hyperfine clock and the Orbitrap axial oscillator are related by the Partition
Lagrangian through the constant $K$, with no additional calibration, no free
parameters, and no residual.

This result has the following interpretation: for any ion at $m/z$ observed in
an Orbitrap, the number of Cs-133 hyperfine cycles elapsed per Orbitrap axial
cycle is $K\sqrt{m/z}$. The mass spectrometer is a clock. Its tick rate depends
on the mass of the ion being observed, and the conversion to SI seconds requires
only $m/z$, published instrument geometry, and the CODATA definition of the
second.

\subsection{Retention Time from Partition Counts}
\label{subsec:rt_result}

Theorem~\ref{thm:rt_from_partition} predicts that for the same ion at two
different retention times, $t_i/t_j = M_{\mathcal{A}}(t_i)/M_{\mathcal{A}}(t_j)$
exactly.

\textbf{Result}: Across all 681 same-$m/z$ pairs (241 ions with multiple
retention times in the dataset):
\begin{equation}
    \max_{\mathrm{pairs}} \left|\frac{M(t_i)/M(t_j) - t_i/t_j}{t_i/t_j}\right|
    = 4.44 \times 10^{-16}
    \label{eq:rt_ratio_result}
\end{equation}
The maximum deviation is floating-point machine epsilon. Retention time ratios
are reproduced exactly by partition count ratios.

More striking: for the isomeric pair $m/z = 230.1387$ (C$_{11}$H$_{19}$NO$_4$)
with elution clusters at RT $\approx 193$ s and RT $\approx 257$ s, the
difference in elution times corresponds to $\Delta n_{\mathrm{cycles}} = 439.0$
additional DDA cycles. The partition count difference between the two clusters
is:
\begin{equation}
    \Delta M = \Delta n_{\mathrm{cycles}} \times M_{\mathrm{cycle}} =
    439.0 \times M_{\mathrm{cycle}}
    \label{eq:isomer_delta_M}
\end{equation}
The 64-second chromatographic separation of these structural isomers is encoded
as 439 additional partition cycles in the acquisition history, recoverable from
partition counts alone without a clock.

\subsection{Composition Inflation Along the Chromatographic Axis}
\label{subsec:inflation_result}

For each ion with $n$ distinct scan events over elution time span $\Delta t$,
the composition inflation count $T(n,2) = 2 \cdot 3^{n-1}$ (Theorem~\ref
{thm:comp_inflation}, $d=2$ for quadrupole + Orbitrap) was computed.

\textbf{Results}:
\begin{itemize}
    \item Maximum depth: $n = 31$ (C$_{15}$H$_{27}$NO$_5$, $m/z = 302.1962$,
          $\Delta t = 557.8$ s)
    \item $T(31,2) = 4.12 \times 10^{14}$ distinguishable partition trajectories
    \item Information density: $7.38 \times 10^{11}$ trajectories per second
    \item Planck depth $n_P = 97$ for this elution
\end{itemize}

The 558-second chromatographic elution of this acylcarnitine encodes
$4.12 \times 10^{14}$ distinguishable acquisition histories in a two-analyser
system. Adding a third analyser channel (IMS, $d=3$) increases this to
$T(31,3) = 3.46 \times 10^{15}$. The exponential growth of distinguishable
trajectories with scan depth is the direct observable consequence of the
composition inflation mechanism derived in Section~\ref{sec:composition}.

\subsection{Summary of Validation Results}
\label{subsec:summary}

\begin{table}[h]
\centering
\caption{Summary of experimental validation results. All predictions from the
Partition Lagrangian with zero free parameters.}
\label{tab:summary}
\begin{tabular}{llll}
\hline
\textbf{Prediction} & \textbf{Predicted} & \textbf{Observed} & \textbf{Error} \\
\hline
$\dot{M}_{\mathrm{Orb}}/\dot{M}_{\mathrm{quad}}$ slope & $-0.5000$ &
$-0.5000$ & $<10^{-15}$ \\
$K$ (Cs/Orb conversion) & $588.016$ & $588.016$ & $0.0000$ ppm \\
RT ratio from partition counts & exact & $4.44\times10^{-16}$ & machine $\epsilon$ \\
Time jump range & 11,200--41,977 & 11,200--41,977 & exact \\
Time jump / transient & $28.3\%$ mean & $28.3\% \pm 2.1\%$ & consistent \\
\hline
\end{tabular}
\end{table}

Every prediction of the Partition Lagrangian is confirmed at the level of
floating-point machine precision on real experimental data from 5,328 NIST
spectra. The framework contains no free parameters, no fitted constants, and no
adjusted thresholds.

% =============================================================================
% SECTION 7: IMPLICATIONS AND DISCUSSION
% =============================================================================

\section{Implications}
\label{sec:implications}

\subsection{The Structure of Reality Requires Interaction}
\label{subsec:structure}

The contact chain self-reinforcement theorem (Theorem~\ref{thm:chain_reinforce})
has a consequence that extends far beyond mechanics: the universe generates
structure not despite interactions but \emph{because} of them. Every contact
event adds a new local negation relationship, which sharpens the boundaries of
every item in the chain.

An item that has been in contact with many other items for a long time is more
defined than an item in isolation. The isolated item exists only by global
negation against all of reality --- the weakest form of existence. The embedded
item exists by a dense network of local negations --- it is not the ground, not
the water, not the organism, not the atmosphere. Each local negation is a
stronger and more specific form of existence than the global negation. The
embedded item is \emph{more real} in the only sense that matters: it is more
precisely distinguished from more things simultaneously.

This gives a partition-theoretic account of why complex systems exist: they are
configurations of high local negation density, where each component's existence
is secured by specific contact relationships with many other components rather
than by diffuse global negation against all of reality. Life, in particular, is
a partition depth chain that maintains itself by continuously generating new
contacts that preserve and deepen the chain's structure.

\subsection{The Partition Second and Metrological Implications}
\label{subsec:metrology}

The SI second is currently defined by the Cs-133 hyperfine transition. This is
a partition count: $\dot{M}_{\mathrm{Cs}} = 9{,}192{,}631{,}770$ Hz. The
partition framework shows that every oscillating system is also a partition
counter, related to the Cs-133 standard by the universal constant $K$ and the
system-specific frequency.

This has the following metrological consequence: the SI second is not special.
It is the agreed reference partition rate. Any oscillating system can serve as a
clock, and the conversion between any two clocks is a ratio of partition rates,
with no free parameters. For mass spectrometry specifically:
\begin{itemize}
    \item An Orbitrap at $m/z = 162.1125$ (carnitine) ticks at
          1,227,842 Hz --- $7{,}487$ times slower than Cs-133
    \item An Orbitrap at $m/z = 650.4$ (large acylcarnitine) ticks at
          613,012 Hz --- $14{,}985$ times slower than Cs-133
    \item The conversion is $K\sqrt{m/z}$ with $K = 588.016$, derived
          from CODATA alone
\end{itemize}
Any LC-MS/MS dataset simultaneously carries multiple independent measurements of
elapsed time --- one per analyser active during acquisition. These measurements
can be cross-validated without a common clock, providing a self-contained
metrological check within each spectrum.

\subsection{Cross-Instrument Calibration}
\label{subsec:cross_calibration}

For two instruments with different analyser types (Orbitrap and FT-ICR)
measuring the same compound at $m/z$:
\begin{equation}
    \frac{\Delta M_{\mathrm{Orb}}}{\dot{M}_{\mathrm{Orb}}(m/z)} =
    \frac{\Delta M_{\mathrm{ICR}}}{\dot{M}_{\mathrm{ICR}}(m/z)} = \Delta t
    \;\mathrm{[s]}
    \label{eq:cross_calibration}
\end{equation}
This equality is a falsifiable prediction. Deviation from equality is a direct
measurement of systematic timing error in either instrument, without any common
external reference. The partition framework provides a basis for cross-instrument
calibration grounded entirely in the physics of the ions being measured.

\subsection{The Four Forces as Partition Depth Gradients}
\label{subsec:forces}

The elimination of forces as separate entities (replacing them with partition
depth minimisation) has a systematic structure. All four fundamental interactions
are partition depth gradients at different scales:

\begin{itemize}
    \item \textbf{Gravity}: Partition depth gradient at macroscopic scales.
          The gradient $\partial\beta_{\mathrm{total}}/\partial D$ drives items
          toward minimum inter-item partition depth. Range: unlimited (because
          every item contributes to every other item's global negation cost).
          Strength: proportional to partition inertia $\mu$ (mass), because
          larger partition count generates stronger depth gradient.

    \item \textbf{Electromagnetism}: Partition orientation gradient. Items with
          anti-parallel boundary orientations ($m_1$ and $m_2$ of opposite sign)
          reduce their total orientation cost by approaching --- local negation is
          more efficient when their boundaries are complementary. Items with
          parallel orientations compete and are repelled. Range: unlimited
          (orientation asymmetry propagates through the contact chain without
          attenuation in vacuum). Strength: proportional to charge magnitude
          (orientation asymmetry $|m|$).

    \item \textbf{Strong force}: Partition depth gradient at nuclear scales.
          Nucleons in contact at partition depth $D = 1$ are in the most
          efficient possible mutual negation state at nuclear scale. The boundary
          at this depth is maintained by the residue parity $s \in
          \{-\frac{1}{2}, +\frac{1}{2}\}$, not by orientation $m$. Two nucleons
          with opposite residue parity ($s_1 = +\frac{1}{2}$, $s_2 =
          -\frac{1}{2}$) form a complementary pair with minimal combined
          boundary thickness. The strong force is short-ranged because the
          residue parity interaction is effective only at $D = 1$ (direct
          contact) --- unlike orientation, parity does not propagate through
          the contact chain.

    \item \textbf{Weak force}: Partition state transition at large $\tau_p$.
          The weak interaction corresponds to a transition between partition
          states with large partition interval --- a slow renegotiation of what
          an item is \emph{not}. The long $\tau_p$ (large $\tau_p$ corresponds
          to small coupling) manifests as the short range of the weak interaction
          (only items in direct contact can undergo parity-changing transitions
          on the timescale of the weak force).
\end{itemize}

\begin{theorem}[Force Hierarchy from Partition Depth]
\label{thm:force_hierarchy}
The relative strengths of the four interactions are determined by the
partition depth at which they operate and the component of the partition
boundary they couple to:
\begin{align}
    \alpha_s &\propto C(\text{depth } 1) = 2 \quad \text{(strong)} \\
    \alpha_{\mathrm{em}} &\propto \frac{|m|}{n} \quad \text{(electromagnetic)} \\
    G_F &\propto \tau_{p,\mathrm{weak}} \quad \text{(weak, large } \tau_p) \\
    G &\propto \lambda_P^2 \quad \text{(gravity, contact geometry)}
\end{align}
The hierarchy strong $\gg$ em $\gg$ weak $\gg$ gravity reflects the hierarchy
of partition depths and boundary components, not four independent coupling
constants.
\end{theorem}

\subsection{Information, Entropy, and the Second Law}
\label{subsec:entropy}

Corollary~\ref{cor:three_forms} establishes that time, entropy, and information
are three expressions of the same residue. The second law of thermodynamics
follows immediately from Theorem~\ref{thm:floor}: since $M$ is strictly monotone
increasing, and entropy is the aggregate of residues whose specific causal links
are not tracked, entropy cannot decrease in any system whose partition count is
increasing --- which is every physical system.

The second law is not a statistical tendency. It is a structural consequence of
the partition framework. It holds not because high-entropy states are
overwhelmingly numerous (though they are) but because the partition count $M$
that generates entropy cannot decrease (Theorem~\ref{thm:floor}).

Shannon's information entropy $H = -\sum_k p_k \log p_k$ is the aggregate of
unresolved causal residues: when we do not track which specific residue $\beta_k$
caused which next partition, we assign probabilities $p_k$ to the possible
causal links and $H$ measures our ignorance of the causal structure. As
causal resolution improves (more specific tracking of which $\beta_k$ causes
which next partition), $H$ decreases. As causal resolution is lost (residues
disperse into the contact chain), $H$ increases. The growth of entropy is the
growth of unresolved causal structure in the contact chain.

\subsection{Virtual Substates and Off-Shell States}
\label{subsec:virtual}

The mean-recovery constraint on virtual substates is a consequence of the
partition framework. For $N$ substates $(s_1, \ldots, s_N)$ of a partition
state $s \in [0,1]$:
\begin{equation}
    \frac{1}{N}\sum_{i=1}^N s_i = s
    \label{eq:mean_recovery}
\end{equation}
Substates with $s_i \notin [0,1]$ are \emph{off-shell virtual substates}. By
the Miracle Principle, for any $s \in [0,1]$, $N \geq 2$, and any $(N-1)$
arbitrary substate values, there exists a unique $s_N = Ns - \sum_{i=1}^{N-1}
s_i$ satisfying mean-recovery.

Off-shell virtual substates are the partition-theoretic description of boundary
states --- states that belong to neither $\mathcal{I}$ nor $\overline{\mathcal{I}}$
but to the boundary between them. The boundary thickness $\beta$ of the
partition is the measure of how far off-shell the boundary states are. Virtual
particles in QFT, chemical transition states, and off-shell virtual substates in
the partition framework are three descriptions of the same object: the
irreducible residue at the boundary of a partition.

\subsection{What the Framework Does Not Claim}
\label{subsec:limits}

The partition framework makes specific, falsifiable predictions. It does not:
\begin{enumerate}
    \item Claim to replace the SI second. The Cs-133 standard remains the agreed
          reference. The framework shows it is one partition rate among many,
          chosen for practical reasons (stability, precision, accessibility).

    \item Claim to improve mass accuracy. The framework is about the metrology
          of time within existing instruments, not about improving $m/z$
          measurement.

    \item Claim that $G$ is derivable. $G$ is identified as irreducible and the
          reason for its irreducibility is explained. Measurement of $G$ remains
          empirical.

    \item Require new hardware. All validation was performed on existing
          instrument types (Q-Orbitrap) using publicly available spectral data.
\end{enumerate}


% =============================================================================
\section{Conclusion}
\label{sec:conclusion}
% =============================================================================

We have shown that the individuation of physical objects --- their existence as
distinct entities --- is a physical process with definite structure. The process
of partitioning an item from the rest of reality requires finite time, requires
the universe to progress during the partition, and generates an irreducible
residue $\beta > 0$ at the boundary. This residue is not waste. It is the causal
propagator of the next partition event: the only reason there is a next process.

From three postulates about this process --- Existence by Negation
(Postulate~\ref{post:negation}), Partition Duration (Postulate~\ref{post:duration}),
and Residue Irreducibility (Postulate~\ref{post:floor}) --- we derived without
additional assumptions:

\begin{enumerate}
\item The partition count $M$ is strictly monotone increasing
      (Theorem~\ref{thm:floor}). This is the origin of the arrow of time, derived
      without entropy, coarse-graining, or boundary conditions.

\item Contact is the minimum viable partition: the most efficient state of
      existence for two distinct items (Theorem~\ref{thm:contact_optimum}).
      Gravity is the observable of partition depth minimisation at macroscopic
      scales (Theorem~\ref{thm:gravity}).

\item Mass, charge, and space emerge from partition structure. Mass is partition
      inertia. Charge is partition boundary orientation. Space is partition depth
      (Definitions~\ref{def:inertia}, \ref{def:charge}, \ref{def:space}).

\item The four standard mass analyser equations (TOF, Orbitrap, FT-ICR,
      quadrupole) are Euler-Lagrange consequences of the Partition Lagrangian
      (Theorem~\ref{thm:four_analysers}).

\item The number of distinguishable partition trajectories grows as $T(n,d) =
      d(d+1)^{n-1}$ (Theorem~\ref{thm:comp_inflation}). Physical time is the
      length of the traversed trajectory; the partition second is the invariant
      time unit (Definition~\ref{def:partition_second}).

\item The gravitational constant $G$ is irreducible to partition counting
      quantities because it encodes the geometry of the contact ground state,
      which is prior to the counting processes that approach it
      (Theorem~\ref{thm:G_irreducible}). It is the one constant in the framework
      that cannot be derived from oscillation rates, tick energies, or branch
      counts.

\item Loschmidt's paradox is a category error. The velocity reversal operation
      requires a frozen universe, which is impossible (Theorem~\ref{thm:frozen}).
      Even if the operation could be performed, the resulting state would have
      decreasing $M$, which Theorem~\ref{thm:floor} forbids. The paradox
      dissolves before any probabilistic argument is required
      (Theorem~\ref{thm:resolution}).
\end{enumerate}

The experimental validation on 5,328 NIST acylcarnitine spectra confirms every
prediction of the framework to machine precision, with zero free parameters
(Table~\ref{tab:summary}).

The contact chain self-reinforcement (Theorem~\ref{thm:chain_reinforce}) gives a
partition-theoretic account of why the universe has structure: every contact
event adds definitional precision to every item in the chain. Structure is the
accumulated residue of contact. Complexity is deep contact chain structure.
Life is a partition depth chain that maintains itself by continuously generating
new contacts that preserve and deepen the chain's structure.

The residue is the causal propagator. Time is its accumulation. Entropy is its
aggregate over unresolved causal links. Information is its transfer between
partition boundaries. These are not three different things described by analogy.
They are the same residue, seen from three different perspectives on the next
partition event.

% =============================================================================
% APPENDICES
% =============================================================================

\appendix

\section{Derivation of Analyser Equations from the Partition Lagrangian}
\label{app:analyser_derivations}

The Partition Lagrangian is:
\begin{equation}
    \mathcal{L}_M = \frac{1}{2}\mu|\dot{x}|^2 + \mu\dot{x} \cdot A_M(x,t)
    - M(x,t)
    \tag{A.1}
\end{equation}
The Euler-Lagrange equation $\frac{d}{dt}\frac{\partial\mathcal{L}_M}{\partial\dot{x}}
- \frac{\partial\mathcal{L}_M}{\partial x} = 0$ gives:
\begin{equation}
    \mu\ddot{x} = -\nabla_x M + \mu(\dot{x} \times (\nabla \times A_M))
    + \mu\frac{\partial A_M}{\partial t}
    \tag{A.2}
\end{equation}

\textbf{TOF}: $M(x) = \frac{qV}{L}x$, $A_M = 0$. Then $\nabla_x M = qV/L$
and $\mu\ddot{x} = -qV/L$. With $\mu = m$ (unit charge), initial velocity 0,
flight path $L$: transit time $T = L\sqrt{m/(2qV)}$. $\checkmark$

\textbf{Orbitrap}: $M(x) = \frac{q\kappa}{2}(x^2 - z^2/2)$ (hyperlogarithmic
potential). Axial equation: $m\ddot{z} = -q\kappa z$. Angular frequency:
$\omega_z = \sqrt{q\kappa/m}$. $\checkmark$

\textbf{FT-ICR}: $A_M = \frac{B}{2}(-y, x, 0)$ (uniform magnetic field
$B\hat{z}$), $M = 0$. Then $\mu\ddot{x} = \mu\dot{x} \times B$. Cyclotron
frequency: $\omega_c = qB/m$. $\checkmark$

\textbf{Quadrupole}: $M(x,y,t) = \frac{U - V\cos(\Omega t)}{r_0^2}(x^2 - y^2)$.
Substituting: $m\ddot{x} = -\frac{2e(U - V\cos\Omega t)}{r_0^2}x$. Setting
$\xi = \Omega t/2$, $a = 4eU/(mr_0^2\Omega^2)$, $q_M = 2eV/(mr_0^2\Omega^2)$:
Mathieu equation $d^2x/d\xi^2 + (a - 2q_M\cos 2\xi)x = 0$. $\checkmark$

\section{Derivation of the Universal Constant K}
\label{app:K_derivation}

\begin{align}
    K &= \frac{\dot{M}_{\Cs} \cdot 2\pi}{\sqrt{e\kappa/u}} \notag \\
      &= \frac{9{,}192{,}631{,}770 \times 2\pi}{\sqrt{1.602176634\times10^{-19}
         \times 10^8 / 1.66053906660\times10^{-27}}} \notag \\
      &= \frac{5.7757 \times 10^{10}}{\sqrt{9.6476 \times 10^{15}}} \notag \\
      &= \frac{5.7757 \times 10^{10}}{9.8224 \times 10^7} \notag \\
      &= 588.016 \;\mathrm{Cs\;cyc\;Orb\;cyc}^{-1}\;\mathrm{u}^{-1/2}
    \tag{B.1}
\end{align}

The conversion $\dot{M}_{\Cs}/\dot{M}_{\Orb}(m/z) = K\sqrt{m/z}$ follows
directly, since $\dot{M}_{\Orb}(m/z) = \sqrt{e(m/z)^{-1}\kappa/u^{-1}}/(2\pi)$
where $(m/z)$ is in atomic mass units and $z$ is the charge state.

\section{The Partition Second in SI Units}
\label{app:partition_second}

For any analyser $\mathcal{A}$ observing an ion at $m/z$ over a duration $\Delta t$:
\begin{align}
    \Delta M_{\mathcal{A}} &= \dot{M}_{\mathcal{A}}(m/z) \cdot \Delta t \tag{C.1}\\
    \Delta t &= \frac{\Delta M_{\mathcal{A}}}{\dot{M}_{\mathcal{A}}(m/z)} \tag{C.2}
\end{align}
In partition seconds:
\begin{equation}
    \Delta t [\mathrm{Ps}] = \frac{\Delta M_{\mathcal{A}}}{\dot{M}_{\mathcal{A}}
    \cdot \dot{M}_{\Cs}} \tag{C.3}
\end{equation}
This equals $\Delta t$ in SI seconds by construction. The partition second is
the SI second expressed entirely in terms of partition counts, without reference
to any specific oscillator. Two instruments with different $\dot{M}$ accumulate
different $\Delta M$ over the same $\Delta t$, but both compute the same
$\Delta t [\mathrm{Ps}]$ via Eq.~(C.2).

% =============================================================================
% BIBLIOGRAPHY
% =============================================================================

\bibliographystyle{plainnat}
\bibliography{references}

\end{document}
